% Statement and cited-input packet. Target proof is separately embargoed.
% Frozen II source SHA256: 74497bfe327fae25c9cdcb17334b3bf42a9793a9df7c0a4404d97568f2f760f0
% Target owner: wbII:thm:candidate-conductor; tag: 03IX

\begin{theorem}[the conductor mechanism for the standard family]
\hypertarget{tag-03IX}{}\label{wbII:thm:candidate-conductor}


\ Assume:
\begin{enumerate}[label=(H\arabic*)]
\item $X$ is a connected compact smooth complex threefold and
$f:X\to B$ is a nonconstant holomorphic map to a smooth compact curve;
\item $W=f^{-1}(p_0)$ is the full scheme-theoretic fibre, is reduced, and
has the normal-crossings local models \eqref{wbII:eq:W:5.3};
\item the normalization $\eta:\widetilde W\to W$ is the degree-six del
Pezzo surface $\mathrm{dP}_6$;
\item $H^2(X;\mathbb Z)=0$.
\end{enumerate}
Then, for every $L\in\operatorname{Pic}(X)$,
\[
H^2(W,(T_X\otimes L)|_W)\ne0
\quad\text{and}\quad
R^2f_*(T_X\otimes L)\ne0.
\tag{W.45}\hypertarget{tag-03HI}{}\label{wbII:eq:W:5.18}
\]
\end{theorem}

% Cited input wbII:eq:W:5.3; statement or labelled formula only.
\[
g=ux,\qquad g=uxy,\qquad\text{or}\qquad g=uxyz
\tag{W.3}\hypertarget{tag-03HS}{}\label{wbII:eq:W:5.3}
\]

% Cited input wbII:lem:curve-flatness; statement or labelled formula only.
\begin{lemma}[automatic properness, flatness, and Cartier fibres]
\hypertarget{tag-03I8}{}\label{wbII:lem:curve-flatness}


\ Let $f:X\to B$ be a nonconstant holomorphic map from a connected
complex manifold to a Riemann surface.  Then $f$ is flat, every holomorphic
vector bundle on $X$ is $f$-flat, and every full scheme-theoretic fibre is an
effective Cartier divisor of pure codimension one.  If $X$ is compact, then
$f$ is proper.
\end{lemma}

% Cited input wbII:prop:analytic-base-change; statement or labelled formula only.
\begin{proposition}[the exact analytic base-change input]
\hypertarget{tag-03II}{}\label{wbII:prop:analytic-base-change}


\ Let $f:X\to B$ be proper between complex spaces, let
$\mathcal F$ be coherent and $f$-flat, and fix $c\in B$.  If
\[
H^3(X_c,\mathcal F|_{X_c})=0,
\tag{W.13}\hypertarget{tag-03ED}{}\label{wbII:eq:W:5.16}
\]
then the canonical fibre map
\[
\beta_c^2:
(R^2f_*\mathcal F)_c\otimes_{\mathcal O_{B,c}}\mathbb C(c)
\longrightarrow H^2(X_c,\mathcal F|_{X_c})
\tag{W.14}\hypertarget{tag-03EE}{}\label{wbII:eq:W:5.17}
\]
is an isomorphism.
\end{proposition}

% Cited input wbII:prop:conductor; statement or labelled formula only.
\begin{theorem}[conductor-supported ambient differential]
\hypertarget{tag-03IK}{}\label{wbII:prop:conductor}

\label{wbII:prop:conductor-differential}
\ In the setting \eqref{wbII:eq:W:5.3}, let $A$ be a line bundle on $S$ with
$\eta^*A\simeq\mathcal O_{\widetilde S}$, and choose a nowhere-vanishing
section $e$ of $\eta^*A$.  There is a unique nonzero
\[
s\in H^0(S,\Omega^1_{\mathcal X}|_S\otimes A)
\tag{W.6}\hypertarget{tag-03HZ}{}\label{wbII:eq:W:5.6}
\]
such that
\[
\eta^*s=\eta^*(dg|_S)\otimes e,\qquad s|_D=0.
\tag{W.7}\hypertarget{tag-03I0}{}\label{wbII:eq:W:5.7}
\]
Its image $\rho(s)$ in $\Omega^1_S\otimes A$ is supported on $D$ and
therefore torsion; its image in
$\widetilde\Omega^1_S\otimes A$ is zero.  Moreover:
\begin{enumerate}[label=(\alph*)]
\item $\rho(s)=0$ if and only if $e=\eta^*\sigma$ for a nowhere-vanishing
$\sigma\in H^0(S,A)$;
\item if $S$ is compact and connected and $A\not\simeq\mathcal O_S$, then
$H^0(S,A)=0$ and $s$ does not come from the conormal term in \CompendiumRef{II}{wbII:eq:W:5.1}.
\end{enumerate}
\end{theorem}

% Cited input wbII:thm:cartier-duality; statement or labelled formula only.
\begin{theorem}[complete Cartier-fibre duality]
\hypertarget{tag-03IY}{}\label{wbII:thm:cartier-duality}


\ Let $X$ be a compact smooth complex threefold,
$i:W\hookrightarrow X$ an effective Cartier divisor, and $E$ a holomorphic
vector bundle.  Put $F=E|_W$ and
\[
\omega_W=(\omega_X\otimes\mathcal O_X(W))|_W.
\]
For $q=0,1,2$ there are natural perfect pairings
\[
H^q(W,F)^\vee
\simeq H^{2-q}(W,F^\vee\otimes\omega_W).
\tag{W.11}\hypertarget{tag-03EB}{}\label{wbII:eq:W:5.11}
\]
Moreover $H^q(W,F)=0$ for $q>2$.  No normality or reducedness hypothesis on
$W$ is used.
\end{theorem}

